\section*{Abstract}
\addcontentsline{toc}{section}{Abstract}

A digitally controlled USB-C power supply was implemented using a charge-pump power-factor-correction (PFC) resonant topology from earlier DTU work and planar magnetics. Burst control was developed for secondary-bulk regulation, startup sequencing, and energy limitation at no load and after load removal.

The gallium-nitride (GaN) half bridge operated at a nominal \SI{1}{\mega\hertz}, with average power controlled by bursts of switching cycles. Feedback samples were supplied through a circular direct memory access (DMA) buffer. Startup precharge was implemented alongside separate burst-duty limits for loaded and no-load operation. The USB-C PD stage provided the tested \SI{5}{\volt}, \SI{9}{\volt}, \SI{15}{\volt}, and \SI{20}{\volt} USB-C settings.

At \SI{240}{\volt}, \SI{50}{\hertz} input, matched measurements at the \SI{20}{\volt} setting yielded \SIrange{73.97}{79.89}{\watt} at the load. Average output voltage was \SIrange{19.973}{19.991}{\volt}, with \SIrange{0.416}{0.452}{\volt} recorded peak-to-peak ripple. Logged power factor was 0.97--0.99 and current total harmonic distortion was \SIrange{2.23}{3.98}{\percent}. Paired input/output readings gave preliminary efficiency estimates of \SIrange{76.5}{77.1}{\percent}, limited by acquisition timing and calibration uncertainty. The tests included no-load operation and load removal. A thermal image labeled \SI{20}{\volt}, \SI{4}{\ampere} showed a \SI{90.4}{\celsius} surface maximum, motivating further loss and thermal characterization.

V2 integrated the MCU, GaN stage, isolated feedback, and USB-C circuit on a four-layer PCB. Lower half-bridge heating was observed after tank and timing changes, while input-bridge heating remained unresolved (\cref{sec:res-v2-testing}). Later V2 tests reached \SI{137.1}{\watt} at the secondary output with voltage droop. Heating in the secondary region and at the input bridge requires further investigation. The V3 edge-card layout retained secondary rectification and omitted the USB-PD output section, giving an estimated net volume reduction of approximately \SI{8}{\percent}. Replacement planar magnetics were not installed. The external regulated USB-PD stage must be included when comparing the volume of complete supplies.

Digitally controlled USB-C operation near \SI{80}{\watt} and higher-power V2 secondary-output operation were demonstrated. Operation at the \SI{100}{\watt} continuous target and \SI{65}{\watt} minimum still needs long-duration thermal and stability tests. Full-range zero-voltage switching, USB-PD conformance, safety, and EMC remain unverified.
